\subsection{Local definition of a Haar measure}

PROPOSITION 9. --- Let G be a locally compact group, V an open subset of G, and \(\mu\) a nonzero positive measure on V having the following property: if U is an open subset of V and if \(s \in G\) is such that \(sU \subset V\) , then the image of the measure \(\mu_{U}\) induced by \(\mu\) on U, under the homeomorphism \(x \mapsto sx\) of U onto sU, is \(\mu_{sU}\) . Then, there exists one and only one left Haar measure \(\alpha\) on G that induces \(\mu\) on V.

For every \(s \in G\) , let \(\mu_{s}\) be the image of \(\mu\) under the homeomorphism \(x \mapsto sx\) of V onto sV. The restriction of \(\mu_{s}\) to \(V \cap sV\) is the image of \(\mu_{s^{-1}V \cap V}\) under the restriction of \(x \mapsto sx\) to \(s^{-1}V \cap V\) ; by hypothesis, this image is \(\mu_{V \cap sV}\) . By translation, one concludes from this that \(\mu_{s}\) and \(\mu_{t}\) have the same restriction to \(sV \cap tV\) for any s, t. Therefore, by Prop. 1 of Ch. III, §2, No.~1, there exists a measure \(\alpha\) on G that induces \(\mu_{s}\) on sV for every s. It is clear that \(\alpha\) is the unique left Haar measure on G inducing \(\mu\) on V.

COROLLARY. --- Let G, G' be two locally compact groups, V (resp. V') an open neighborhood of the neutral element of G (resp. G'), and \(\varphi\) a local isomorphism of G' with G (GT, III, §1, No.~3, Def. 2) defined on V', such that \(\varphi(V') = V\) . Let \(\alpha'\) be a left Haar measure on G', and \(\alpha'_{V'}\) its restriction to \(V'\) . Then \(\varphi(\alpha_{V'}')\) is the restriction to V of a unique left Haar measure \(\alpha\) on G.

Let \(V_1\) be an open neighborhood of \(e\) in \(G\) such that \(V_1V_1^{-1} \subset V\) . Let \(\mu\) be the restriction of \(\varphi(\alpha_{V'}')\) to \(V_1\) . Let \(U\) be an open subset of \(V_1\) and let \(s \in G\) be such that \(sU \subset V_1\) . Then \(s \in V_1V_1^{-1} \subset V\) , therefore \(s = \varphi(s')\) for some \(s' \in V'\) . Let \(x \in U\) . Then \(x = \varphi(x')\) for some \(x' \in V'\) , therefore \(sx = \varphi(s')\varphi(x') = \varphi(s'x')\) since \(sx \in sU \subset V\) . Since the left translations in \(G'\) preserve \(\alpha'\) , one sees that \(V_1\) and \(\mu\) satisfy the conditions of Prop. 9. Let \(\alpha\) be the left Haar measure on \(G\) inducing \(\mu\) on \(V_1\) . For every \(t \in V\) , there exists an open neighborhood \(W\) of \(e\) in \(V_1\) such that \(tW \subset V\) . Then the restriction of \(\varphi(\alpha_{V'}')\) to \(tW\) may be deduced by translation from the restriction of \(\mu\) to \(W\) , hence is the restriction of \(\alpha\) to \(tW\) . Therefore \(\varphi(\alpha_{V'}')\) is the restriction of \(\alpha\) to \(V\) .

One says that \(\alpha\) is deduced from \(\alpha'\) by means of the local isomorphism \(\varphi\) .

Example. --- The Haar measure on T obtained in No.~2, Example 3 may be deduced from the Lebesgue measure on R by a local isomorphism of R with T.

